\documentclass[11pt]{article}
\usepackage[a4paper,margin=25mm]{geometry}
\usepackage{amsmath,amssymb,amsthm,bm}
\usepackage{booktabs}
\usepackage{xcolor}
\usepackage{microtype}
\usepackage[hidelinks]{hyperref}

\newcommand{\E}{\mathbb E}
\newcommand{\Var}{\operatorname{Var}}
\newcommand{\mc}{\mu_{\mathrm c}}
\newtheorem{proposition}{Proposition}
\newtheorem{corollary}{Corollary}
\theoremstyle{remark}
\newtheorem{remark}{Remark}

\title{The Size--Variance Exponent on Heterogeneous Graphs\\
\large A decomposition of $\beta$ for the relative GLV}
\author{}
\date{30 September 2026}

\begin{document}
\maketitle

\begin{center}
\fcolorbox{black!35}{blue!4}{\begin{minipage}{0.91\textwidth}
\textbf{Main result.}
At leading order in $1/C$ the exponent is the HDMFT's, and analytically
\[
  \beta_\infty\simeq\beta_{\rm ad}^\Phi-\varepsilon_\chi-\varepsilon_k .
\]
Here $\beta_{\rm ad}^\Phi$ is the adiabatic size--volatility slope of a single firm, normalised by
its realised fitness amplitude.  It equals $1$
when firms fully track their fitness, and falls when the field is rough.  $\varepsilon_\chi$ is
a small fluctuation-selection term ($0.02$--$0.03$).  $\varepsilon_k$ is the degree channel: how
strongly larger surviving firms sit on higher degrees.  It follows from an explicit survivor
ensemble built on three justified facts:
\begin{enumerate}
  \item a firm's mean size equals its mean fitness (exact);
  \item that mean fitness is Gaussian, with amplitude proportional to $\sqrt\kappa$;
  \item a firm survives the observation window if its mean fitness exceeds its deepest dip, an
        extreme-value law $P=\Phi(\bar\Phi/f)^{T/2\tau}$.
\end{enumerate}
With order parameters from the HDMFT and no fitting, it reproduces the HDMFT's degree channel
and size spread to within $0.015$ for $\alpha\ge2.5$.  Quenched neighbourhoods add $1/C$
corrections ($\varepsilon_n$, $\varepsilon_z$).  With every input from the HDMFT, the result is
within $0.07$ of simulation at $N=8000$ for $\alpha\ge2.5$.  The theory predicts, and simulation confirms, that $\beta$
depends on the observation window.  The same ensemble gives the multiscaling: the higher degree
moments at fixed size grow faster than the first, because the smallest survivors are pinned at the
minimum degree, and the HDMFT reproduces the simulated ratio curve $\zeta_q/\zeta_1$ to within
$0.05$ for $\alpha=2.5$ and $3.5$ (Section~\ref{sec:multiscaling}).  Two things are not closed: $\beta_{\rm ad}$ on light tails
(a nonlinear one-firm effect), and $\alpha\approx2.2$, where the field is carried by about four
neighbours and the Gaussian closure fails.
\end{minipage}}
\end{center}

\section{Setup}

We use the notation of the HDMFT appendix: $\mc=-\mu>0$ is competition,
$\kappa=k/C$ is the degree relative to the mean, and $S_i=x_i/m$ is the relative size.  At
zero reciprocity the effective process of a firm in degree class $\kappa$ is
\begin{equation}
  \frac{d\ln S}{dt}=\Phi(t)-S+\lambda\left(S^{-1}-1\right),
  \qquad
  \Phi(t)=1-g(t)-\mc B_\kappa(t)+\sigma\eta_\kappa(t),
  \label{eq:process}
\end{equation}
with $\E[\eta_\kappa(t)\eta_\kappa(t')]=q_\kappa(t,t')$.  The fitness $\Phi$ excludes the
firm's own self-limitation.  For an uncorrelated kernel, $P_{pr}=v_pv_r/c$, the closures of the
appendix factorise:
\begin{equation}
  B_\kappa=\kappa\bar B,\qquad q_\kappa(t,t')=\kappa\,Q(t,t'),
  \label{eq:factor}
\end{equation}
with $\bar B$ and $Q$ neighbour averages that are the same for every class.  The structural
cutoff breaks \eqref{eq:factor} only for the hubs.  So the random part of the field is
$\sigma\sqrt\kappa\,\xi(t)$, where $\xi$ is a class-independent Gaussian process with covariance
$Q$.

\paragraph{Observables.}
We fix a horizon $\Delta$ ($=0.5$ in the production runs) and set $G_i=\ln S_i(t+\Delta)-\ln
S_i(t)$.  The volatility is $v_i$, the MAD estimate of the standard deviation of $G_i$ over the
late window, and $\bar S_i$ is the time-averaged size.  The local exponent is
\begin{equation}
  \beta(S)=-\frac{d\ln\E[v\mid\bar S=S]}{d\ln S},
\end{equation}
and the reported $\beta$ is its least-squares average over the decline branch.  In a finite fit
window $\beta$ is the window average of $\beta(S)$, and every statement below should be read
that way.

\section{Fixed degree: the adiabatic volatility law}

\begin{proposition}[Adiabatic volatility law]\label{prop:adiabatic}
Linearise \eqref{eq:process} about the firm's mean size, where $\bar S=\bar\Phi$ at
$\lambda\to0$.  The log-size responds to the fitness through a low-pass filter,
\begin{equation}
  \delta\widehat{\ln S}(\omega)=\frac{\delta\widehat\Phi(\omega)}{\bar S+i\omega},
  \label{eq:filter}
\end{equation}
so the volatility of a firm of class $\kappa$ and size $S$ is
\begin{equation}
  v^2=\sigma^2\kappa\,A^2(S),
  \qquad
  A^2(S)=\int_{-\infty}^{\infty}\frac{d\omega}{2\pi}\;
     \frac{4\sin^2(\omega\Delta/2)\,\widehat Q(\omega)}{S^2+\omega^2},
  \label{eq:A}
\end{equation}
where $\widehat Q$ is the spectrum of the fluctuating part of $\xi$.  The function $A$ does not
depend on $\kappa$.
\end{proposition}

\begin{proof}
Write $S=\bar S(1+u)$ with $u$ small.  Then $\dot u=-\bar S u+\delta\Phi$ to first order, which
is \eqref{eq:filter}.  Next, $G=u(t+\Delta)-u(t)$ carries the factor $|e^{i\omega\Delta}-1|^2=
4\sin^2(\omega\Delta/2)$.  Finally, $\delta\Phi=\sigma\sqrt\kappa\,\delta\xi$ by
\eqref{eq:factor}.
\end{proof}

There are three regimes, set by the correlation rate $\omega_c$ of the field and by the horizon.
\begin{enumerate}
  \item $S\Delta\gg1$ and $S\gg\omega_c$ (adiabatic). The firm tracks its fitness,
        $\ln S\approx\ln\Phi$, so
        \begin{equation}
          v=\frac{\sqrt{D}}{S},\qquad D=\Var[\Phi(t+\Delta)-\Phi(t)],
          \label{eq:adiabatic}
        \end{equation}
        and $\beta_A=1$.  This limit does not need linearity in $\delta\Phi$, only the
        separation of time scales, so it holds firm by firm with $D_i$ the realised increment
        variance of that firm's own field.
  \item $\omega_c\ll S\ll\Delta^{-1}$, with a rough field ($\widehat Q\sim\omega^{-2}$ up to
        $\Delta^{-1}$).  Then $A^2\propto\int d\omega\,\Delta^2\omega^{-2}\omega^2/(S^2+
        \omega^2)\propto S^{-1}$, so $\beta_A=\tfrac12$.
  \item $S\ll\omega_c$.  The firm integrates the field without relaxing, $A$ is
        independent of $S$, and $\beta_A=0$.  This is the small-firm plateau, and the
        immigration floor sits inside it.
\end{enumerate}

\begin{corollary}[Homogeneous graphs]\label{cor:homog}
On a regular or fully connected graph $\kappa\equiv1$, so $\beta=\beta_A$.  As the fit window
moves into regime~1, $\beta\to1$.
\end{corollary}
This matches the simulations: $\beta=1.02$ on the fully connected graph and $0.99$ on the
regular graph at $N=32000$.  At $\alpha=6$, where the degrees are nearly homogeneous, the
realised-amplitude slope $\beta_D$ of Table~\ref{tab:check} is $1.03$.  The corollary also
explains the horizon dependence in the numerics appendix: raising $\Delta$ moves a fixed size
window from regime~2 towards regime~1, so $\beta_A$ and hence $\beta$ grow with the horizon.

\paragraph{Check of the degree independence of $A$.}
At $N=16000$ on the $\alpha=2.5$ graph, $A=v S/\sqrt k$ binned in size agrees across the four
degree quartiles to within a few per cent at every size.  A joint regression
$\ln v=c-a\ln S+b\ln k$ gives $b=0.50\pm0.03$ across economies (column $b_k$ of
Table~\ref{tab:check}: $0.42$--$0.54$ for every $\alpha$).  The $\sqrt\kappa$ prefactor of
\eqref{eq:A} is therefore quantitative.  The local slope of $A$ rises from $0.4$ at the peak
($S\approx1$) to $0.8$ at $S\approx8$, which is the crossover from regime~2 to regime~1.  A
one-parameter Ornstein--Uhlenbeck spectrum in \eqref{eq:A} reproduces its window slope ($0.666$
against $0.672$).

\section{Mixing degrees: a sum rule}

\begin{proposition}[Sum rule]\label{prop:sum}
Assume $v=\sigma\sqrt\kappa\,A(S)\,\zeta$ with $\zeta$ independent of $\kappa$ given $S$.  Then,
exactly,
\begin{equation}
  \beta(S)=\beta_A(S)-\varepsilon(S),
  \qquad
  \beta_A=-\frac{d\ln A}{d\ln S},
  \qquad
  \varepsilon=\frac{d\ln\E[\sqrt\kappa\mid S]}{d\ln S}.
  \label{eq:sum}
\end{equation}
\end{proposition}

\begin{proof}
Take $\E[v\mid S]=\sigma A(S)\,\E[\zeta]\,\E[\sqrt\kappa\mid S]$ and differentiate its logarithm.
\end{proof}

A graph lowers $\beta$ only to the extent that larger firms sit, on average, on larger degrees.
The size of the degree heterogeneity does not matter on its own; what matters is how strongly
size selects on it.  The same argument holds with the realised amplitude $\sqrt{D_i}$ of
\eqref{eq:adiabatic} in place of $\sqrt\kappa$, which gives $\beta=\beta_D-\varepsilon_D$.

\subsection{Intermediate case: a two-class graph}

Take two degree classes $\kappa_1<\kappa_2$ with weights $\rho_{1,2}$.  Assume that size is
proportional to the noise amplitude (Proposition~\ref{prop:size} below):
$\ln S=\tfrac12\ln\kappa+\eta$, with $\eta\sim\mathcal N(m,V_\eta)$ independent of the class.
Write $y=\ln S$, $c_j=\tfrac12\ln\kappa_j$ and $\delta=c_2-c_1$.  The posterior weight of the
upper class is logistic in $y$,
\begin{equation}
  \operatorname{logit}\pi_2(y)=\ln\frac{\rho_2}{\rho_1}
     +\frac{\delta}{V_\eta}\left(y-m-\frac{c_1+c_2}{2}\right),
\end{equation}
and $\E[\sqrt\kappa\mid y]=e^{c_1}\bigl[1+(e^\delta-1)\pi_2(y)\bigr]$.  Hence
\begin{equation}
  \varepsilon(y)=\frac{(e^\delta-1)\,\tfrac{\delta}{V_\eta}\,\pi_2(1-\pi_2)}
                      {1+(e^\delta-1)\pi_2}.
\end{equation}
Three facts follow.
\begin{enumerate}
  \item The dip in $\beta(S)$ is localised.  It has width $\sim V_\eta/\delta$ in $\ln S$ and
        depth $\simeq\delta^2/4V_\eta$ for small $\delta$.  Outside it, $\beta=\beta_A$.
  \item The sum rule has a fixed budget:
        \begin{equation}
          \int_{-\infty}^{\infty}\bigl[\beta_A-\beta\bigr]\,d\ln S
          =\ln\frac{\E[\sqrt\kappa\mid S\to\infty]}{\E[\sqrt\kappa\mid S\to0]}
          =\tfrac12\ln\frac{\kappa_2}{\kappa_1}.
          \label{eq:budget}
        \end{equation}
  \item A fit window of log-width $L$ that contains the whole dip returns
        $\beta_{\rm fit}\simeq\beta_A-\tfrac12\ln(\kappa_2/\kappa_1)/L$.
\end{enumerate}

For any degree distribution the budget \eqref{eq:budget} is
$\tfrac12\ln(\kappa_{\max}/\kappa_{\min})$ at most.  Under the fixed dilution $C=N/40$ the range
$\kappa_{\max}/\kappa_{\min}\le(N/C)/\kappa_{\min}$ does not grow with $N$.  The degree channel
therefore cannot send $\beta\to0$ as $N\to\infty$, and a finite $\beta_\infty$ is the expected
outcome.

\subsection{Power-law graph: an empirical shortcut}

This subsection closes the degree channel with two survivor statistics, the survival exponent
$s$ and the size spread $V_S$, taken from simulation.  Section~\ref{sec:ensemble} derives the
same channel from the dynamics, with no simulation input.

\begin{proposition}[Size tracks amplitude]\label{prop:size}
Among surviving firms, the size is proportional to the noise amplitude:
$d\ln\bar S/d\ln\sqrt\kappa=1$.  Equivalently $\bar S\propto k^{\delta}$ with $\delta=\tfrac12$.
\end{proposition}
\begin{proof}[Derivation]
Average \eqref{eq:process} over a stationary window.  The left side vanishes, so
\begin{equation}
  \bar S=\bar\Phi+\lambda\bigl(\overline{S^{-1}}-1\bigr)\simeq\bar\Phi
  \label{eq:stationarity}
\end{equation}
for every firm above the floor.  The time-averaged size \emph{is} the time-averaged fitness
(numerically, $\mathrm{corr}(\bar S,\bar\Phi)=1.000$ and $\bar S/\bar\Phi=1.000$ on every graph).
Split $\bar\Phi=r_\kappa+\sigma a_s x$.  Here $r_\kappa=1-\bar g-\mc\kappa\bar B$ is the
deterministic load, and $\sigma a_s x$ is the static random field
$-(\sigma/\sqrt C)\sum_jA_{ij}z_{ij}\bar S_j$.  This field has conditional standard deviation
$\sigma a_s$, with $a_s^2=\sum_jA_{ij}\bar S_j^2/C$, and standardised value $x$.  On the decline
branch the median load is small ($|r_\kappa|\lesssim0.1$ for $\alpha\le2.5$), so
$\bar S\simeq\sigma a_s x$.  The size is then a static amplitude times an $O(1)$ z-score.  Since
$a_s\propto\sqrt\kappa$ to leading order, $\delta=\tfrac12$.
\end{proof}
The derivation needs $|r_\kappa|\ll\bar S$.  Truncating a Gaussian with a large load
$\mc\kappa\bar B$ would give $\delta\approx0.3$, and this is where the light-tailed graphs
depart slightly.  The forward elasticity $f_k$ in Table~\ref{tab:check} is $0.92$--$1.05$ for
every $\alpha$, and the thesis value is $\delta=0.56\pm0.02$.

Now close \eqref{eq:sum} with a log-Gaussian joint law, $\ln\bar S=\tfrac12\ln\kappa+\eta$ with
$\eta\perp\kappa$.  The regression of $\tfrac12\ln\kappa$ on $\ln S$ is then linear with
constant slope,
\begin{equation}
  \varepsilon=\frac{\Var(\tfrac12\ln\kappa)}{\Var\ln S}=\frac{V_\kappa}{4V_S}.
  \label{eq:eps-closure}
\end{equation}
Survival thins the degree tail.  The surviving degrees are Pareto with exponent $\alpha+s$, with
$s=0.82$--$0.99$ on all five graphs, so $\ln(\kappa/\kappa_{\min})\sim\mathrm{Exp}(\alpha-1+s)$
and
\begin{equation}
  V_\kappa=\frac{1}{(\alpha-1+s)^2},
  \qquad
  \boxed{\;\varepsilon_k=\frac{1}{4V_S\,(\alpha-1+s)^2}\;}
  \label{eq:eps-alpha}
\end{equation}
The measured $V_\kappa$ matches $1/(\alpha-1+s)^2$ to within $0.013$.  The log-variance of the
surviving sizes, $V_S\simeq0.19$--$0.24$, is nearly topology independent.  It is set by the
disorder, since the degree part $V_\kappa/4$ is at most a quarter of it.  With $V_S\approx0.22$ and
$s\approx0.9$, \eqref{eq:eps-alpha} gives $\varepsilon_k\approx1.14/(\alpha-0.1)^2$.  The degree
channel therefore lowers $\beta$ by at most $\approx0.25$ even at $\alpha=2.2$.

\section{Numerical check}

\begin{table}[h]
\centering
\small\setlength{\tabcolsep}{4pt}
\begin{tabular}{lcccc|ccccc|ccc}
\toprule
$\alpha$ & $\beta$ & $\beta_A$ & $\varepsilon_k$ & $\beta_A-\varepsilon_k$
 & $\varepsilon_k^{\rm th}$ & $b_k$ & $f_k$ & $s$ & $V_\kappa$ / th
 & $\beta_D$ & $\varepsilon_D$ & $V_R$\\
\midrule
2.2 & 0.222 & 0.440 & 0.229 & 0.211 & 0.244 & 0.42 & 0.94 & 0.87 & 0.221 / 0.234 & 0.74 & 0.55 & 0.128\\
2.5 & 0.414 & 0.581 & 0.174 & 0.407 & 0.183 & 0.50 & 1.05 & 0.99 & 0.151 / 0.162 & 0.77 & 0.36 & 0.057\\
3.0 & 0.595 & 0.715 & 0.149 & 0.565 & 0.148 & 0.50 & 1.00 & 0.85 & 0.117 / 0.123 & 0.89 & 0.33 & 0.040\\
3.5 & 0.689 & 0.788 & 0.102 & 0.686 & 0.112 & 0.53 & 0.92 & 0.94 & 0.081 / 0.084 & 0.93 & 0.25 & 0.033\\
6.0 & 0.894 & 0.916 & 0.034 & 0.882 & 0.033 & 0.54 & 0.95 & 0.82 & 0.030 / 0.030 & 1.03 & 0.16 & 0.035\\
\bottomrule
\end{tabular}
\caption{Medians over economies, $N=8000$, $C=200$, $\mu=-1.76$, $\sigma=1.75$,
$\lambda=10^{-3}$, eight economies per graph ($\alpha=2.2$: seven, one freezes), each fitted on
its own decline branch.  $\varepsilon_k^{\rm th}$ is \eqref{eq:eps-alpha} evaluated with the
measured $V_S$ and $s$.  The $V_\kappa$ column gives the measured value and $1/(\alpha-1+s)^2$.
$\beta_D$ and $\varepsilon_D$ repeat the sum rule with the realised field amplitude $\sqrt{D_i}$,
and $V_R=\Var\ln\sqrt{D_i/k_i}$.  Source: \texttt{scripts/beta\_theory.py}.}
\label{tab:check}
\end{table}

Three things can be read off the table.
\begin{itemize}
  \item The sum rule \eqref{eq:sum} holds to within $0.04$ on every graph.  The residual is
        the gap between the median in $\beta$ and the mean in $\varepsilon$.
  \item The closure \eqref{eq:eps-alpha} predicts the degree channel from $\alpha$, $s$ and
        $V_S$ to within $0.02$.
  \item Most of the topology dependence is in $\beta_A$: $1-\beta_A=0.56$ against
        $\varepsilon_k=0.23$ at $\alpha=2.2$.  The fixed-degree slope is itself lowered by
        heterogeneity.
\end{itemize}

\section{The survivor ensemble: an analytic degree channel}
\label{sec:ensemble}

The shortcut \eqref{eq:eps-alpha} needs $s$ and $V_S$ from simulation.  Both are properties of
which firms survive the observation window.  This section derives them at leading order in $1/C$,
where the Gaussian cavity field of the HDMFT is exact.  Every ingredient below was checked inside
the HDMFT itself, with paths whose fitness was reconstructed exactly by inverting
\eqref{eq:process}.

\paragraph{Ingredients.}
\begin{enumerate}
  \item \emph{Size is mean fitness}, \eqref{eq:stationarity}: $\bar S=\bar\Phi$ for every
        surviving firm (median difference below $0.01$).
  \item \emph{The mean fitness is Gaussian at fixed degree.}  Over a window of length $T$,
        \begin{equation}
          \bar\Phi=r_\kappa+a_\kappa x,\qquad x\sim\mathcal N(0,1),\qquad
          a_\kappa=\sigma\sqrt{\kappa\,q_T},\qquad r_\kappa=1-\bar g-\mc\kappa\bar B .
        \end{equation}
        Here $q_T=T^{-2}\!\int\!\!\int_TQ$ is the variance of the window-averaged field.  Around
        its window mean, the fitness fluctuates with amplitude $f_\kappa=\sigma\sqrt{\kappa(q_0-q_T)}$.
        In simulation both amplitudes scale as $\kappa^{0.50}$, and $x$ is Gaussian to within an
        excess kurtosis of $0.2$--$0.5$ for $\alpha\ge2.5$.
  \item \emph{Survival is an extreme-value event.}  A firm survives the window if its mean
        fitness exceeds its deepest fitness dip.  Otherwise it spends long enough at negative
        fitness to decay to the floor.  For a stationary Gaussian fluctuation with correlation
        time $\tau$, the window contains $n\simeq T/2\tau$ roughly independent blocks, so
        \begin{equation}
          P_{\rm surv}(u)=\Phi_{\mathcal N}(u)^{\,n},\qquad u=\frac{\bar\Phi}{f_\kappa},
          \qquad n=\frac{T}{2\tau}.
          \label{eq:ev}
        \end{equation}
\end{enumerate}

Two simpler survival laws fail, and the way they fail is informative.  A sharp level crossing
(Rice rate) gives a threshold proportional to $f_\kappa$.  A drawdown law (death when $\int\Phi$
falls by $\ln(S/\text{floor})$) gives a threshold proportional to $f_\kappa^2$.  In the HDMFT the
survival curve depends on $u$ alone and is the same at low and high degree:

\begin{center}
\small
\begin{tabular}{lcccccccc}
\toprule
$u=\bar\Phi/f_\kappa$ & 0 & 0.25 & 0.5 & 0.75 & 1.0 & 1.25 & 1.5 & 1.75\\
\midrule
low-degree half  & 0.00 & 0.01 & 0.07 & 0.25 & 0.38 & 0.61 & 0.84 & 0.92\\
high-degree half & 0.00 & 0.01 & 0.07 & 0.21 & 0.32 & 0.55 & 0.77 & 0.82\\
\bottomrule
\end{tabular}
\end{center}

This is the scale-free level-crossing form.  Its soft shape is that of \eqref{eq:ev} with
$n\approx9$ ($\alpha=3.5$, $T=60$, $\tau=3.4$ gives $n=8.9$).  A hard threshold at $u\approx1.2$
has the right location but the wrong shape.  It would miss the survivors just below threshold,
and those set the low tail of the size distribution.  In simulation the threshold rises slightly
faster than $f_\kappa$ ($\kappa^{0.74-0.90}$ against $\kappa^{0.57}$).  That is a finite-$C$
effect, since within-class amplitude heterogeneity is absent from the HDMFT.

\paragraph{The degree channel.}
The joint law of degree and size among survivors is explicit:
\begin{equation}
  p(\kappa,S)\;\propto\;p(\kappa)\;\frac1{a_\kappa}\,
  \phi\!\left(\frac{S-r_\kappa}{a_\kappa}\right)\;
  \Phi_{\mathcal N}\!\left(\frac{S}{f_\kappa}\right)^{T/2\tau},
  \label{eq:joint}
\end{equation}
and $\varepsilon_k(S)=d\ln\E[\sqrt\kappa\mid S]/d\ln S$ follows by one quadrature over $\kappa$
at each $S$.  Differentiating \eqref{eq:joint} gives the exact identity
\begin{equation}
  \varepsilon_k(S)=S\;\mathrm{Cov}_{\kappa\mid S}\!\left(\tfrac12\ln\kappa,\;
     -\frac{S-r_\kappa}{a_\kappa^2}+\frac{n}{f_\kappa}\,
     \frac{\phi(S/f_\kappa)}{\Phi_{\mathcal N}(S/f_\kappa)}\right).
  \label{eq:eps-cov}
\end{equation}
Two forces act on the degree of a firm of given size.  The first term is the Gaussian cost of
reaching $S$.  It is cheaper for large $a_\kappa\propto\sqrt\kappa$, so large firms are drawn
towards high degree ($\varepsilon_k>0$), and the load $-\mc\kappa\bar B$ in $r_\kappa$ pushes the
other way.  The second term is the survival pressure, which acts near the threshold.  The same
law explains $\delta\approx\tfrac12$: the threshold is in units of $f_\kappa\propto\sqrt\kappa$,
so surviving sizes scale close to $\sqrt\kappa$ whatever the load.  The model gives
$\delta=0.40$ at $T=60$, against $0.43$--$0.51$ in the HDMFT and $0.47$--$0.52$ in simulation.
The load term pulls it slightly below $\tfrac12$.

The inputs are the HDMFT order parameters $\bar g$, $\bar B$, $q_0$, $q_T$ and the correlation
time $\tau$, together with the degree distribution and the observation window $T$.  None is
fitted to $\beta$.

\begin{table}[h]
\centering
\small\setlength{\tabcolsep}{4pt}
\begin{tabular}{lcc|cc|cc|cc}
\toprule
$\alpha$ & $\tau$ & $n$ & surv HD & surv \eqref{eq:ev} & $\varepsilon_k$ HD & $\varepsilon_k$
 \eqref{eq:joint} & $V_S$ HD & $V_S$ \eqref{eq:joint}\\
\midrule
2.2 & 2.41 & 12.4 & 0.139 & 0.094 & 0.287 & 0.332 & 0.149 & 0.140\\
2.5 & 2.73 & 11.0 & 0.124 & 0.096 & 0.220 & 0.235 & 0.157 & 0.156\\
3.5 & 3.38 & 8.9 & 0.148 & 0.133 & 0.111 & 0.115 & 0.170 & 0.174\\
6.0 & 3.62 & 8.3 & 0.175 & 0.152 & 0.036 & 0.035 & 0.174 & 0.180\\
\bottomrule
\end{tabular}
\caption{The survivor ensemble \eqref{eq:joint} with the extreme-value survival law, against
the HDMFT's own survivors.  Order parameters are read from the HDMFT solution, and nothing is
fitted.  The model runs $10$--$30\%$ low on the survival fraction, but the degree channel and
the size spread are reproduced to within $0.015$ for $\alpha\ge2.5$.  Replacing \eqref{eq:ev}
by a hard threshold doubles $\varepsilon_k$ (to $0.41$ at $\alpha=2.5$), because it narrows
$V_S$ to $0.09$.  Source: \texttt{scripts/beta\_survivor\_theory.py}.}
\label{tab:ensemble}
\end{table}

\section{Closing the within-degree channel}
\label{sec:within}

\subsection{The field amplitude is a neighbour sum}

Let $d_j=\Var[S_j(t+\Delta)-S_j(t)]$ be the increment variance of firm $j$.  At zero reciprocity
the increments of different neighbours are nearly uncorrelated, so the field increment variance
of firm $i$ is
\begin{equation}
  D_i\simeq\frac{\sigma^2}{C}\sum_jA_{ij}z_{ij}^2\,d_j .
  \label{eq:nbrsum}
\end{equation}
The measured ratio of $D_i$ to the right side is $0.95$--$0.99$ on every graph.  What matters is
how few neighbours carry the sum.  The effective number,
$n_{\rm eff}=(\sum_jc_j)^2/\sum_jc_j^2$ with $c_j=z_{ij}^2d_j$, is $24$ at $\alpha=6$ and only
$4$ at $\alpha=2.2$.  On a heavy-tailed graph a firm's field is carried by a handful of large
neighbours.

Two quenched sums over the same neighbours therefore enter the dynamics.  The static one,
$a_s^2=\sum_jA_{ij}\bar S_j^2/C$, sets the size through \eqref{eq:stationarity}.  The dynamic
one, \eqref{eq:nbrsum}, sets the volatility through \eqref{eq:adiabatic}.  Selection on size
reaches the volatility in two ways.

\subsection{Neighbour composition, $\varepsilon_n$}

Averaging over the couplings gives $D_i^{(0)}=(\sigma^2/C)\sum_jA_{ij}d_j$.  A firm whose
neighbours are large has both a large $a_s$ (so it is big) and a large $D^{(0)}$ (so it is
noisy).  With neighbours drawn with weight $w_j\propto k_j$ (configuration model) and the delta
method, the covariance of the two log-sums per neighbour is set by the cross-moment
\begin{equation}
  c_{ds}=\frac{\E_w[d\,\bar S^2]}{\E_w[d]\,\E_w[\bar S^2]},
  \qquad
  \mathrm{Cov}\bigl(\ln\sqrt{D^{(0)}/\kappa},\,\ln\bar S\bigr)\simeq\frac{c_{ds}-1}{4k}.
\end{equation}
With $\ln\bar S=\ln a_s+\ln x$ and $x$ independent of the neighbours, the reverse regression gives
\begin{equation}
  \boxed{\;\varepsilon_n=\frac{c_{ds}-1}{4V_S}\,\E\!\left[k^{-1}\right].\;}
  \label{eq:eps-n}
\end{equation}
Here $c_{ds}$ measures how much the neighbours that carry the static amplitude also carry the
dynamic one.  It grows with the degree tail, from $3.4$ at $\alpha=6$ to $6.3$ at $\alpha=2.2$.

\subsection{Selection on the couplings, $\varepsilon_z$}

At fixed neighbours, a firm becomes big by drawing favourable couplings $z_{ij}$ on its big
neighbours.  These same $|z_{ij}|$ enter \eqref{eq:nbrsum} squared.  The $z_{ij}$ are Gaussian,
so conditioning them on the static z-score $x$ is exact:
\begin{equation}
  \E[z_{ij}\mid x]=\frac{\bar S_j\,x}{\sqrt{\sum_lA_{il}\bar S_l^2}},\qquad
  \Var[z_{ij}\mid x]=1-\frac{\bar S_j^2}{\sum_lA_{il}\bar S_l^2}.
\end{equation}
Substituting into \eqref{eq:nbrsum} gives
\begin{equation}
  \E[D_i\mid x]=D_i^{(0)}\bigl[1+\theta_i(x^2-1)\bigr],
  \qquad
  \theta_i=\frac{\sum_jA_{ij}d_j\bar S_j^2}{\sum_jA_{ij}d_j\,\sum_jA_{ij}\bar S_j^2}
  \simeq\frac{c_{ds}}{k_i}.
  \label{eq:zsel}
\end{equation}
The concentration $\theta_i$ is of order $1/n_{\rm eff}$.  When the field is spread over many
neighbours, conditioning on one linear combination of the couplings barely moves their squares.
When a few neighbours dominate, it does.  By \eqref{eq:stationarity}, $x=(\bar S-r_\kappa)/
(\sigma a_s)$ is fixed by the size (measured correlation $1.000$), and
\begin{equation}
  \boxed{\;\varepsilon_z=\frac12\,\frac{d\ln\E\bigl[1+\theta(x^2-1)\mid S\bigr]}{d\ln S}
  \;\approx\;\frac{\theta x^2}{1+\theta(x^2-1)}\;
  \Bigl(1-\frac{V_{a_s}}{V_S}\Bigr).\;}
  \label{eq:eps-z}
\end{equation}
The last factor is $d\ln x/d\ln S$.  It is unity at fixed $a_s$, and is reduced by the part of
the size variance that comes from the amplitude itself.  The z-selection saturates at
$\varepsilon_z\to1$ for $\theta x^2\gg1$.  This is the regime approached by a firm fed by a
single neighbour.

\subsection{Result}

Collecting the channels,
\begin{equation}
  \boxed{\;\beta=\beta_{\rm ad}^{\rm nb}-\varepsilon_k-\varepsilon_n-\varepsilon_z,\;}
  \label{eq:full}
\end{equation}
with $\varepsilon_k$ from \eqref{eq:eps-alpha}, $\varepsilon_n$ from \eqref{eq:eps-n} and
$\varepsilon_z$ from \eqref{eq:eps-z}.  Here $\beta_{\rm ad}^{\rm nb}$ is the adiabatic slope with the neighbour-sum amplitude \eqref{eq:nbrsum} (Section~\ref{sec:bad}).  Apart from $\beta_{\rm ad}^{\rm nb}$, every input is a
neighbour-weighted moment of single-firm statistics ($d$, $\bar S$, $k$) or a size-level
quantity ($V_S$, $x$).  No quenched information about a firm's actual neighbourhood enters, so
all inputs are outputs of the degree-resolved effective process.

\begin{table}[h]
\centering
\small\setlength{\tabcolsep}{3.5pt}
\begin{tabular}{lccc|cc|cc|cc|cc|c}
\toprule
$\alpha$ & $n_{\rm eff}$ & $c_{ds}$ & $\E[k^{-1}]$ & $\varepsilon_k$ & th
 & $\varepsilon_n$ & th & $\varepsilon_z$ & th & $\beta_{\rm ad}^{\rm nb}$ & $\beta$ & pred\\
\midrule
2.2 & 4.0 & 6.28 & 0.020 & 0.229 & 0.230 & 0.105 & 0.106 & 0.161 & 0.172 & 0.729 & 0.209 & 0.188\\
2.5 & 8.8 & 4.61 & 0.011 & 0.174 & 0.169 & 0.050 & 0.048 & 0.108 & 0.076 & 0.736 & 0.414 & 0.459\\
3.5 & 17.4 & 3.77 & 0.006 & 0.102 & 0.109 & 0.027 & 0.024 & 0.064 & 0.036 & 0.888 & 0.691 & 0.720\\
6.0 & 23.5 & 3.35 & 0.005 & 0.034 & 0.035 & 0.016 & 0.014 & 0.040 & 0.025 & 0.965 & 0.894 & 0.891\\
\bottomrule
\end{tabular}
\caption{The three amplitude channels, measured as slopes of $\E[\sqrt{\,\cdot\,}\mid S]$ on
the estimator's decline bins and predicted by \eqref{eq:eps-alpha}, \eqref{eq:eps-n} and
\eqref{eq:eps-z}.  Same ensemble as Table~\ref{tab:check}, with $\Delta=0.5$ on the $0.25$
output grid.  ``pred'' is \eqref{eq:full} with the three closed forms and the measured
$\beta_{\rm ad}^{\rm nb}$.  Source: \texttt{scripts/beta\_within\_degree.py}.}
\label{tab:within}
\end{table}

The degree and composition channels are predicted to within $0.007$.  The coupling-selection
channel is predicted to within $0.03$.  It runs low by $0.015$--$0.03$ on the lighter tails,
where the neglected off-diagonal terms of \eqref{eq:nbrsum} are comparable to the effect.  The
full prediction \eqref{eq:full} is within $0.05$ of the measured $\beta$ on every graph.  The
decomposition also shows where heterogeneity acts.  At $\alpha=2.2$ the drop $1-\beta\approx0.8$
splits into $0.27$ from non-adiabaticity, $0.23$ from degree, $0.11$ from neighbour composition
and $0.16$ from coupling selection.  The last two are the within-degree channel, and both scale
as $c_{ds}/k$, the inverse effective number of neighbours.

\paragraph{Consequences.}
\begin{itemize}
  \item \emph{System size.}  Both within-degree channels carry a factor $\E[k^{-1}]$, which
        halves when $N$ doubles at $C=N/40$.  They shrink more slowly than that.  From $N=8000$
        to $16000$, $\varepsilon_n+\varepsilon_z$ falls from $0.27$ to $0.24$ at $\alpha=2.2$ and
        from $0.16$ to $0.12$ at $\alpha=2.5$, factors of $0.90$ and $0.73$ rather than $0.5$.  Two
        things slow the decay.  $c_{ds}$ grows with the heaviest neighbour weights ($4.6\to5.5$ at
        $\alpha=2.5$).  And on the decline branch $\E[k^{-1}]$ itself falls by less than half
        ($0.020\to0.013$ at $\alpha=2.2$), because the survivors shift towards higher degrees.  At
        the same time $\varepsilon_k$ grows towards its HDMFT value ($0.174\to0.213$ at
        $\alpha=2.5$, against $0.214$), so the degree channel also carries a finite-$C$
        correction, of the same size as the fall in $\varepsilon_n+\varepsilon_z$.  The
        finite-size drift of $\beta$ is therefore a trade between channels, and \eqref{eq:full}
        tracks it: predicted $0.180$ and $0.455$ against measured $0.152$ and $0.431$ at
        $N=16000$.  Over the measured range the approach to the limit is not linear in $1/N$.
        A large-$N$ law needs the $N$-dependence of $c_{ds}$, $V_S$ and the survivors' degree
        mix, which we have not derived.  At fixed $C$, by contrast, $\E[k^{-1}]$ does not
        fall while $c_{ds}$ keeps growing with $\kappa_{\max}=N/C$.  This is consistent with the
        observed fall of $\beta$ towards zero.
  \item \emph{Homogeneous graphs.}  Here $c_{ds}\to1+O(\text{size dispersion})$ and
        $\E[k^{-1}]=1/C$, so $\varepsilon_n,\varepsilon_z=O(1/C)$ and $\beta\to\beta_{\rm ad}\to1$.
  \item \emph{What remains.}  The one input not yet closed here is $\beta_{\rm ad}^{\rm nb}$, the filter
        slope of Proposition~\ref{prop:adiabatic} with the neighbour-sum amplitude.
        Section~\ref{sec:bad} closes it with the two-time HDMFT.
\end{itemize}

\section{Closing $\beta_{\rm ad}$ with the two-time HDMFT}
\label{sec:bad}

\paragraph{Two adiabatic slopes.}
$\beta_{\rm ad}=1-d\ln\E[vS/\sqrt D\mid S]/d\ln S$ depends on which amplitude $D$ normalises the
volatility, and two are used below.  $\beta_{\rm ad}^{\rm nb}$ uses the neighbour sum
$D_i^{\rm nb}=\sum_jz_{ij}^2d_j$ of \eqref{eq:nbrsum}.  This is the amplitude against which the
channels $\varepsilon_n$, $\varepsilon_z$ of Section~\ref{sec:within} are defined, and the
$\beta_{\rm ad}$ of Table~\ref{tab:within}.  At $C\to\infty$ it becomes the class amplitude of
the HDMFT.  $\beta_{\rm ad}^\Phi$ uses the realised increment variance of the full fitness,
$D_i^\Phi=\Var[\Phi_i(t+\Delta)-\Phi_i(t)]$ with $\Phi_i=1-g-(\alpha S)_i$, as in
\eqref{eq:adiabatic}.  By Proposition~\ref{prop:sum} they differ by a selection channel,
\begin{equation}
  \beta_{\rm ad}^\Phi-\beta_{\rm ad}^{\rm nb}=\varepsilon_r,
  \qquad
  \varepsilon_r=\frac{d\ln\E\bigl[\sqrt{D^\Phi/D^{\rm nb}}\mid S\bigr]}{d\ln S},
\end{equation}
which collects the mean-coupling part of the field and the off-diagonal terms dropped in
\eqref{eq:nbrsum}.  It is the finite-$C$ counterpart of the HDMFT selection term
$\varepsilon_\chi$ below.  Measured on the same eight economies per graph, $\varepsilon_r=-0.04$,
$0.01$, $0.04$, $0.05$ ($\alpha=2.2$, $2.5$, $3.5$, $6$), and
$\beta=\beta_{\rm ad}^\Phi-\varepsilon_k-\varepsilon_n-\varepsilon_z-\varepsilon_r$ closes to
within $0.03$ on every graph (\texttt{scripts/beta\_ad\_estimators.py}).  Comparisons with the
HDMFT must pair like with like: class amplitude with $\beta_{\rm ad}^{\rm nb}$, realised
amplitude with $\beta_{\rm ad}^\Phi$.

The same script integrates each economy twice, at the two tolerances used by the scripts of this
note ($10^{-4}$ and $10^{-6}$).  A chaotic trajectory does not survive a change of tolerance, so
the late window is resampled, and the eight-economy medians of $\beta$ and $\beta_{\rm ad}$ move by
up to $0.05$ and $0.07$.  This is the resampling noise found in the thesis's tolerance check
(per-economy shifts up to $0.04$), not a bias.  It sets the resolution of every comparison in
this note at about $0.03$--$0.05$.  All simulation values below come from the $10^{-4}$ runs,
the same economies as Table~\ref{tab:within}.

\paragraph{The single-firm replay is exact.}
Drive \eqref{eq:process} with a firm's own measured fitness $\Phi_i(t)$ and the full
nonlinearity, starting from its true size.  The replay reproduces $\beta_{\rm ad}^\Phi$ to within
$0.005$ on every graph, so $\beta_{\rm ad}$ is a property of the one-firm process and the law of
its field.  Linearising the replay with a constant rate $\bar S_i$, as in \eqref{eq:A}, loses
accuracy.  It is right on heavy tails ($0.67$ against $0.65$ at $\alpha=2.2$) but runs
$0.07$--$0.11$ low on light tails, where it cannot exceed $1$ while the measurement reaches
$1.01$.  The shortfall is a genuine nonlinear effect and not an estimator artefact: the
standard-deviation estimate of the slope is higher still.  The linear filter therefore gives the
mechanism but not the number.

\paragraph{One field spectrum per graph suffices.}
Replacing every firm's own field spectrum by the average spectrum of the decline branch changes
the linear-response $\beta_{\rm ad}$ by at most $0.012$.  The shape of the field spectrum is not
firm-specific.  What varies across graphs is its roughness, the ratio of the field-increment
variance to the field variance.  This ratio is $0.036$ at $\alpha=6$ and $0.43$ at $\alpha=2.2$,
and it is a single-site statistic.  In the diagonal, annealed approximation of
\eqref{eq:nbrsum},
\begin{equation}
  \frac{\Var[\Delta h]}{\Var[h]}\simeq\frac{\E_w[d_j]}{\E_w[\Var_tS_j]},
  \qquad w_j\propto k_j,
  \label{eq:rough}
\end{equation}
which gives $0.408$, $0.073$, $0.056$ and $0.030$ against measured values of $0.429$, $0.088$,
$0.068$ and $0.036$ ($\alpha=2.2$, $2.5$, $3.5$, $6$).  Heavy tails roughen the field because
neighbours are drawn in proportion to degree.  Surviving hubs are large
(Proposition~\ref{prop:size}), and large firms relax fast, at rate $\bar S_j$, so their
fluctuations carry the high frequencies.  A rough field pushes the firms at the start of the
decline branch, where $\bar S\Delta\approx0.5$, out of the adiabatic regime, and $\beta_{\rm ad}$
falls.

\paragraph{The self-consistent class field reproduces $\beta_{\rm ad}$.}
Because $\beta_{\rm ad}$ depends only on the one-firm process and a graph-level field law, the
degree-resolved two-time HDMFT should capture it.  Its classes run \eqref{eq:process} in a
self-consistent Gaussian field.  We solve it with 25 classes from the realised degree
sequence, the SCM kernel and 200 paths per class
(\texttt{hdmft.solve\_\allowbreak twotime\_\allowbreak heterogeneous}), and measure the
slope of $vS/\sqrt{D_p}$ on its surviving paths, with $D_p$ the class field-increment variance
from $q_p$.

\begin{table}[h]
\centering
\small\setlength{\tabcolsep}{4pt}
\begin{tabular}{l|cc|cc|ccc|cc}
\toprule
 & \multicolumn{2}{c|}{class / neighbour sum} & \multicolumn{2}{c|}{realised $\Phi$} & & & & & \\
$\alpha$ & $\beta_{\rm ad}^{\rm HD}$ & $\beta_{\rm ad}^{\rm nb}$
 & $\beta_{\rm ad}^{\rm HD,\Phi}$ & $\beta_{\rm ad}^\Phi$
 & $\beta^{\rm HD}$ & $\varepsilon_k^{\rm HD}$ & $\beta_{\rm ad}^{\rm HD}-\varepsilon_k^{\rm HD}$
 & $\beta_{\rm th}$ & $\beta^{\rm sim}$\\
\midrule
2.2 & 0.643 & 0.729 & 0.680 & 0.697 & 0.397 & 0.268 & 0.375 & ---   & 0.209\\
2.5 & 0.753 & 0.736 & 0.782 & 0.749 & 0.566 & 0.214 & 0.539 & 0.388 & 0.414\\
3.5 & 0.902 & 0.888 & 0.937 & 0.918 & 0.826 & 0.109 & 0.793 & 0.759 & 0.691\\
6.0 & 0.966 & 0.965 & 0.987 & 1.014 & 0.945 & 0.035 & 0.931 & 0.900 & 0.894\\
\bottomrule
\end{tabular}
\caption{Two-time HDMFT against simulation ($N=8000$ graphs, $\mu=-1.76$, $\sigma=1.75$,
$\lambda=10^{-3}$).  Each simulated $\beta_{\rm ad}$ is paired with its HDMFT counterpart: the
class amplitude with the neighbour sum $D^{\rm nb}$, the realised path fitness with
$D^\Phi$.  Simulation columns are medians over the eight economies of Table~\ref{tab:within}
(\texttt{scripts/beta\_ad\_estimators.py}).  $\beta_{\rm th}=\beta^{\rm HD}-\varepsilon_n-\varepsilon_z$ with every
input of \eqref{eq:eps-n} and \eqref{eq:eps-z} ($c_{ds}$, $\E[k^{-1}]$, $V_S$, $x$) taken from
HDMFT paths.  $V_S$ adds the static composition spread
$\tfrac14\E[k^{-1}](\E_w[\bar S^4]/\E_w[\bar S^2]^2-1)$, which the class-level field does not
carry.  Source: \texttt{scripts/beta\_ad\_hdmft.py}.}
\label{tab:hdmft}
\end{table}

Table~\ref{tab:hdmft} gives three results.
\begin{enumerate}
  \item Paired like with like, $\beta_{\rm ad}$ is predicted by the HDMFT to within $0.02$
        (class against neighbour sum) and $0.03$ (realised against realised) for
        $\alpha\ge2.5$.  At $\alpha=2.2$ the realised pair still agrees ($0.680$ against
        $0.697$), but the class pair is $0.09$ apart, since the HDMFT is not converged there.
        The HDMFT field roughness does not match the simulation in detail
        ($0.135$ against $0.088$ at $\alpha=2.5$).  The agreement in $\beta_{\rm ad}$ is
        therefore an agreement of the full one-firm process, not of a single roughness number.
  \item The HDMFT's own exponent obeys $\beta^{\rm HD}\simeq\beta_{\rm ad}^{\rm HD}-
        \varepsilon_k^{\rm HD}$ to within $0.03$.  This is expected: its class field has no
        quenched neighbourhoods, so it carries only the adiabatic and degree channels.  The
        gap $\beta^{\rm HD}-\beta^{\rm sim}$ ($0.19$, $0.15$, $0.14$, $0.05$) is the
        within-degree channel of Section~\ref{sec:within}.
  \item Taking the within-degree inputs from the HDMFT paths too gives a $\beta$ with no
        simulation input: $0.388$, $0.759$ and $0.900$ against $0.414$, $0.691$ and $0.894$
        for $\alpha=2.5$, $3.5$ and $6$.
\end{enumerate}

\paragraph{What sets $\beta_{\rm ad}$.}
Two further checks inside the HDMFT narrow this down.  First, normalising by the class-level
$D_p$ mixes in a selection effect: near-threshold survivors are the ones whose realised
fluctuations happened to be small.  With each path's realised increment variance instead, the
adiabatic slope is $0.680$, $0.782$, $0.937$ and $0.987$ ($\alpha=2.2$, $2.5$, $3.5$, $6$), and
the selection term is only $\varepsilon_\chi=0.02$--$0.03$.  These realised values match the
simulation's $\beta_{\rm ad}^\Phi$ ($0.697$, $0.749$, $0.918$, $1.014$) to within $0.03$.  Second, linear response with the class
spectrum gives a flat $0.66$--$0.71$ and misses the trend.  Linear response on the survivors'
own fitness paths recovers it on heavy tails ($0.688$, $0.746$ against $0.680$, $0.782$), but runs
$0.09$--$0.12$ low on light tails.  Survivors see a slightly rougher field than their class
($D/\Var=0.163$ against $0.143$ at $\alpha=2.5$), because survival removes paths with deep slow
dips.  The remaining light-tail gap is a second-order nonlinear response that we have not
derived.  $\beta_{\rm ad}$ is therefore a property of the one-firm process.  It is computable from
a one-dimensional stochastic equation, and linear response gets it analytically only on heavy
tails.

\paragraph{Where the Gaussian closure fails.}
At $\alpha=2.2$ the HDMFT fixed point does not converge (relative error $15$--$18\%$ even at
twice the iterations, and $g$ keeps drifting).  Its survivors also sit on too-low degrees
($\E[k^{-1}]=0.043$ against $0.020$ in simulation), so the within-degree inputs come out wrong
and $\beta_{\rm th}$ is meaningless there.  This is the regime where
$n_{\rm eff}\approx4$.  A field carried by four neighbours is not Gaussian, so the assumption
behind the cavity step fails.  For $\alpha=2.2$ the theory therefore stands on
Table~\ref{tab:within}: closed forms with simulation inputs, and $\beta_{\rm ad}$ from either
source.  A sparse (few-neighbour, non-Gaussian) cavity field is what the heaviest tails would
need.

\section{Assembled prediction and a new test}
\label{sec:assembled}

At leading order in $1/C$ the exponent is the HDMFT's own, $\beta_\infty\equiv\beta^{\rm HD}$,
since the class field carries no quenched neighbourhoods.  The survivor ensemble gives it
analytically,
\begin{equation}
  \beta_\infty\simeq\beta_{\rm ad}^{\rm HD,\Phi}-\varepsilon_\chi-\varepsilon_k,
  \label{eq:beta-inf}
\end{equation}
with $\varepsilon_k$ from \eqref{eq:joint}, and $\beta_{\rm ad}^{\rm HD,\Phi}$, $\varepsilon_\chi$ from the
one-firm process.  The within-degree channels of Section~\ref{sec:within} are the first
quenched correction, $\beta_N=\beta_\infty-\varepsilon_n-\varepsilon_z$, both proportional to
$c_{ds}\E[k^{-1}]$.

\begin{center}
\small
\begin{tabular}{lcccc}
\toprule
$\alpha$ & $\beta_\infty$, \eqref{eq:beta-inf} & $\beta_\infty=\beta^{\rm HD}$
 & $\beta_{8000}$ (theory) & $\beta_{8000}$ (sim)\\
\midrule
2.2 & 0.320 & 0.397 & ---   & 0.209\\
2.5 & 0.519 & 0.566 & 0.388 & 0.414\\
3.5 & 0.797 & 0.826 & 0.759 & 0.691\\
6.0 & 0.933 & 0.945 & 0.900 & 0.894\\
\bottomrule
\end{tabular}
\end{center}

$\beta_{8000}$ (theory) is $\beta^{\rm HD}-\varepsilon_n-\varepsilon_z$ with every input taken
from HDMFT paths, the $\beta_{\rm th}$ of Table~\ref{tab:hdmft}; it has no simulation input.
The analytic form \eqref{eq:beta-inf} runs $0.01$--$0.05$ below $\beta^{\rm HD}$, because the
ensemble overestimates $\varepsilon_k$ on the heavier tails (Table~\ref{tab:ensemble}) and the
split $\beta^{\rm HD}\simeq\beta_{\rm ad}^{\rm HD}-\varepsilon_k^{\rm HD}$ holds only to $0.03$.
The theory is within $0.03$ of simulation at $\alpha=2.5$ and $6$, and $0.07$ above it at
$\alpha=3.5$.  The thesis's $1/N$ extrapolations ($0.79$ at $\alpha=3.5$, and a loosely
constrained $0.47$ at $\alpha=2.5$) lie $0.04$ and $0.10$ below $\beta^{\rm HD}$.  They assume a
$1/N$ approach, which Section~\ref{sec:within} shows is not established, so they are not a sharp
test.  At $\alpha=2.2$ the $1/C$ expansion is not controlled ($n_{\rm eff}\approx4$) and the HDMFT
does not converge, so no theoretical $\beta_{8000}$ is given.

\paragraph{Window dependence.}
Through $n=T/2\tau$ and $q_T$, \eqref{eq:joint} predicts that $\beta$ depends on the
observation window $T$, which the thesis never varied.  A longer window is a stricter survival
test: fewer firms survive, survivors are selected further out in the size tail, and
$\varepsilon_k$ grows.  Simulation confirms the direction:

\begin{center}
\small
\begin{tabular}{lc|cc|cc|c}
\toprule
$\alpha$ & $T$ & surv sim & surv \eqref{eq:ev} & $\varepsilon_k$ sim & $\varepsilon_k$
 \eqref{eq:joint} & $\beta$ sim\\
\midrule
2.5 & 30  & 0.271 & 0.179 & 0.132 & 0.185 & 0.474\\
2.5 & 60  & 0.227 & 0.097 & 0.179 & 0.234 & 0.468\\
2.5 & 120 & 0.110 & 0.050 & 0.201 & 0.335 & 0.390\\
3.5 & 30  & 0.219 & 0.212 & 0.074 & 0.085 & 0.689\\
3.5 & 60  & 0.147 & 0.132 & 0.095 & 0.113 & 0.708\\
3.5 & 120 & 0.097 & 0.080 & 0.129 & 0.168 & 0.703\\
\bottomrule
\end{tabular}
\end{center}

The trends agree.  The levels differ because the simulation is at finite $C$, where survival
is higher at $\alpha=2.5$.  The size of the shift is overestimated at $T=120$, where $q_T$ is
extrapolated as $1/T$ from the $T=60$ HDMFT window.  At $\alpha=3.5$ the rise of
$\varepsilon_k$ with $T$ is offset by $\beta_{\rm ad}$, and $\beta$ is flat.  At $\alpha=2.5$
$\beta$ falls by $0.08$ between $T=60$ and $120$.  Any comparison with the empirical
$\beta\approx0.15$--$0.20$ therefore has to match the empirical panel length as well as the
horizon $\Delta$.

\section{Multiscaling from the survivor ensemble}
\label{sec:multiscaling}

The higher moments of the size-conditioned volatility define exponents
$\E[v^q\mid\bar S=S]\sim S^{-\zeta_q}$, with $\zeta_1=\beta$.  A volatility law of fixed shape
gives $\zeta_q/\zeta_1=q$.  The data give $1,2.0,2.6,2.9$ for $q=1,\dots,4$, and the thesis shows
that the model bends below the line only on heterogeneous graphs.  This section derives that
bend from the survivor ensemble \eqref{eq:joint}.

\subsection{Splitting the moments}

Write $v=\sqrt\kappa\,A(S)\,\zeta$ as in Proposition~\ref{prop:sum}.  If $\zeta$ and $\kappa$ are
independent given $S$, the $q$-th moment factorises and
\begin{equation}
  \zeta_q=\zeta_q^A-e_q,
  \qquad
  e_q=\frac{d\ln\E[\kappa^{q/2}\mid S]}{d\ln S},
  \label{eq:zeta-split}
\end{equation}
where $\zeta_q^A$ is the $q$-th exponent of the degree-normalised volatility $v/\sqrt\kappa$, and
$e_1=\varepsilon_k$.  Dividing by $\zeta_1$,
\begin{equation}
  q-\frac{\zeta_q}{\zeta_1}
  =\frac{(e_q-q\,e_1)+(q\,\zeta_1^A-\zeta_q^A)}{\zeta_1}.
  \label{eq:bend}
\end{equation}
The bend below the fixed-shape line has two sources: a degree moment that grows faster than the
$q$-th power of the first, and a degree-normalised volatility that is not of fixed shape.  Both
are divided by $\zeta_1=\beta$, so a small exponent amplifies whatever bend there is.

\begin{proposition}[Multiscaling criterion]\label{prop:multiscaling}
Expand the cumulant generating function of $\ln\kappa$ given $S$,
$\ln\E[\kappa^{q/2}\mid S]=\tfrac q2 m(S)+\tfrac{q^2}{8}s^2(S)+O(q^3)$.  Then
\begin{equation}
  e_q-q\,e_1=\frac{q(q-1)}{8}\,\frac{ds^2}{d\ln S}+O(q^3),
  \label{eq:criterion}
\end{equation}
so the degree channel bends the ratio curve below $q$ if and only if the conditional spread
$s^2=\Var(\ln\kappa\mid S)$ grows with size.
\end{proposition}

\begin{proof}
Differentiate the expansion in $\ln S$: $e_q=\tfrac q2m'+\tfrac{q^2}{8}(s^2)'$, and subtract
$q\,e_1=\tfrac q2m'+\tfrac q8(s^2)'$.
\end{proof}

A log-Gaussian closure with $\eta\perp\kappa$, as in \eqref{eq:eps-closure}, has constant
$s^2$ and hence no multiscaling.  The survivor ensemble does not.

\subsection{Why the conditional spread grows}

Take the load small ($|r_\kappa|\ll S$, as on the decline branch for $\alpha\le2.5$) and use
$a_\kappa=a_1\sqrt\kappa$, $f_\kappa=f_1\sqrt\kappa$.  For a Pareto degree law,
$p(\kappa)\propto\kappa^{-\alpha}$ on $[\kappa_{\min},\kappa_{\max}]$, change variable in
\eqref{eq:joint} to $u=S/\sqrt\kappa$.  The conditional law at fixed size is
\begin{equation}
  p(u\mid S)\;\propto\;u^{2\alpha-2}\,\phi\!\left(\frac u{a_1}\right)
  \Phi_{\mathcal N}\!\left(\frac u{f_1}\right)^{n},
  \qquad
  \frac S{\sqrt{\kappa_{\max}}}<u<\frac S{\sqrt{\kappa_{\min}}},
  \label{eq:u-law}
\end{equation}
and $\E[\kappa^{q/2}\mid S]=S^q\,\E[u^{-q}\mid S]$.  The size enters only through the edges of
the support.  Three consequences follow.
\begin{enumerate}
  \item Without degree cutoffs, $\ln\kappa=2\ln S-2\ln u$ with $u$ independent of $S$.  Then
        $e_q=q$ exactly, $s^2=4\Var\ln u$ is constant, and there is no multiscaling: size is
        pure amplitude.
  \item The lower cutoff pins the smallest survivors to $\kappa_{\min}$, where $s^2=0$.  As $S$
        grows the edge $S/\sqrt{\kappa_{\min}}$ moves out through the bulk of \eqref{eq:u-law},
        the range of admissible degrees opens, and $s^2$ rises towards $4\Var\ln u$.  By
        Proposition~\ref{prop:multiscaling}, this rise is the multiscaling.
  \item The upper cutoff carries almost no weight.  For $\kappa\gg S^2/f_1^2$ the survival factor
        tends to $2^{-n}$ ($n\approx8$--$12$), so hubs of a given size do not survive and the
        edge $S/\sqrt{\kappa_{\max}}$ sits where the density has already vanished.
\end{enumerate}
The degree multiscaling is therefore a crossover: it lives on the range of sizes over which the
lower cutoff releases, and it fades once $s^2$ has saturated.  This fits the reading of
Moran et al.\ (2024), who treat the empirical sublinearity as a subleading correction rather than an
asymptotic law.

\subsection{Numerical check}

We compute every term of \eqref{eq:bend} three ways, each on its own decline branch with the
thesis estimator (bin means of $y^q$, per economy): on the HDMFT's surviving paths (leading order
in $1/C$), from the ensemble \eqref{eq:joint} with HDMFT order parameters interpolated in $\kappa$
and drawn on the realised degree sequence (nothing fitted), and in simulation.

\begin{table}[h]
\centering
\small\setlength{\tabcolsep}{4pt}
\begin{tabular}{l|ccc|cc|cc}
\toprule
 & \multicolumn{3}{c|}{$e_4/e_1$} & \multicolumn{2}{c|}{$\zeta_4^A/\zeta_1^A$}
 & \multicolumn{2}{c}{$e_1/\zeta_1$}\\
$\alpha$ & HD & ensemble & sim & HD & sim & HD & sim\\
\midrule
2.2 & 4.88 & 5.59 & 5.45 & 3.90 & 3.68 & 0.80 & 1.25\\
2.5 & 5.29 & 6.15 & 5.69 & 3.87 & 3.73 & 0.42 & 0.48\\
3.5 & 5.92 & 6.38 & 5.40 & 3.96 & 3.80 & 0.15 & 0.16\\
6.0 & 5.47 & 5.40 & 5.15 & 3.98 & 3.87 & 0.04 & 0.04\\
\bottomrule
\end{tabular}
\caption{The two sources of the bend in \eqref{eq:bend} at $q=4$.  The fixed-shape value of both
ratios is $4$.  The degree moments are strongly superlinear, with ratios $e_q/e_1\approx
1,\,2.2,\,3.7,\,5.4$ for $q=1,\dots,4$ on every graph.  In the HDMFT the degree-normalised
volatility keeps its shape.  The bend amplitude is set by $e_1/\zeta_1=\varepsilon_k/\beta$.
Same graphs and operating point as Table~\ref{tab:check}; eight economies per graph ($\alpha=2.2$:
seven).  Source: \texttt{scripts/zeta\_survivor\_theory.py}.}
\label{tab:zeta-split}
\end{table}

\begin{table}[h]
\centering
\small\setlength{\tabcolsep}{4pt}
\begin{tabular}{l|c|c|c}
\toprule
$\alpha$ & HDMFT & $\zeta_q^{A,\rm HD}-e_q^{\rm ens}$ & simulation\\
\midrule
2.2 & 1.79, 2.40, 2.87 & 1.80, 2.28, 2.19 & 1.86, 2.51, 3.01\\
2.5 & 1.86, 2.60, 3.23 & 1.85, 2.48, 2.79 & 1.88, 2.63, 3.20\\
3.5 & 1.93, 2.79, 3.55 & 1.95, 2.84, 3.61 & 1.94, 2.80, 3.60\\
6.0 & 1.98, 2.94, 3.88 & 1.99, 2.97, 3.93 & 1.96, 2.87, 3.73\\
\midrule
data & \multicolumn{3}{c}{2.0, 2.6, 2.9}\\
\bottomrule
\end{tabular}
\caption{The ratio curve $\zeta_q/\zeta_1$ for $q=2,3,4$.  The middle column is the theory: the
HDMFT's degree-normalised exponents combined with the degree channel of the survivor ensemble
through \eqref{eq:zeta-split}.  Source: \texttt{scripts/zeta\_survivor\_theory.py}.}
\label{tab:zeta-ratio}
\end{table}

Four things can be read off the tables.
\begin{enumerate}
  \item At leading order the bend is all degree.  In the HDMFT $\zeta_q^A/\zeta_1^A$ stays within
        $0.13$ of $q$ at $q=4$, while $e_q/e_1$ exceeds $q$ by $0.9$--$1.9$.  The second term of
        \eqref{eq:bend} is a small correction, and the mechanism is Proposition~\ref{prop:multiscaling}.
  \item The HDMFT reproduces the simulated ratio curve to within $0.05$ for $\alpha=2.5$ and
        $3.5$, and to within $0.15$ at $\alpha=2.2$ and $6$.  The ensemble, with no simulation
        input, tracks the HDMFT through $q=3$ (to within $0.12$ at $\alpha=2.5$ and $0.05$
        above) and is within $0.15$ of simulation there.  At $q=4$ on the heavy tails it
        overshoots $e_4$, because the fourth moment is carried by the top degree classes, whose
        order parameters are interpolated from coarse HDMFT bins.
  \item The ratio curve is a leading-order property and $\beta$ is not.  The HDMFT is off in
        $\zeta_1$ by $0.14$ at $\alpha=2.5$ ($0.525$ against $0.389$) and still gets the ratios
        right.  This is consistent with the thesis finding that the ratios are converged in $N$
        while $\beta$ drifts: the within-degree channels of Section~\ref{sec:within} move
        $\zeta_1$ at order $1/C$.  We have not shown why they leave the ratios nearly unchanged.
  \item The bend grows with the tail for two reasons, visible in \eqref{eq:bend}: $e_1=
        \varepsilon_k$ rises and $\zeta_1=\beta$ falls.  Their ratio goes from $0.04$ at
        $\alpha=6$ to $0.8$--$1.25$ at $\alpha=2.2$.
\end{enumerate}

\paragraph{Limits.}
In simulation the degree-normalised volatility is not of fixed shape ($\zeta_4^A/\zeta_1^A=
3.68$--$3.87$), so part of the simulated bend comes from the within-degree amplitude $R_i$ rather
than from degree.  Simulation and HDMFT reach nearly the same curve by slightly different routes.
The factorisation \eqref{eq:zeta-split} holds to $0.01$ at $q=1$ and to about $0.1$ at $q=4$ in
the HDMFT.  In simulation it fails at $\alpha=2.2$ ($0.31$ at $q=4$), the same few-neighbour
regime where the Gaussian cavity fails.  Finally, \eqref{eq:criterion} is a second-cumulant
effect, so it bends the ratio curve roughly as $q(q-1)$.  It cannot produce the hard saturation
of the data after $q=3$, which is why the model's curve bends more gently than the data's.

\section*{Summary}
\begin{enumerate}
  \item At fixed degree a firm tracks its fitness, so $v=\sqrt D/S$ and $\beta\to1$ on
        homogeneous graphs (Proposition~\ref{prop:adiabatic}).
  \item Exact sum rule: a graph lowers $\beta$ only by the amount size selects on noise
        amplitude (Proposition~\ref{prop:sum}).
  \item Leading order in $1/C$: the survivor ensemble \eqref{eq:joint} gives $\varepsilon_k$ with
        no fitting.  Its ingredients are size $=$ mean fitness, a Gaussian static field, and
        extreme-value survival.  It reproduces the HDMFT's degree channel and size spread to
        within $0.015$ ($\alpha\ge2.5$).
  \item First correction in $1/C$: quenched neighbourhoods, $\varepsilon_n+\varepsilon_z$, both
        proportional to $c_{ds}\E[k^{-1}]$.  Because $c_{ds}$ grows with $N$ and
        $\varepsilon_k$ converges to its HDMFT value from below, $\beta(N)$ approaches its limit
        more slowly than $1/N$ over the measured range.
  \item Prediction confirmed in simulation: $\beta$ and $\varepsilon_k$ depend on the observation
        window $T$.
  \item Multiscaling: the degree channel bends $\zeta_q/\zeta_1$ below $q$ if and only if
        $\Var(\ln\kappa\mid S)$ grows with size (Proposition~\ref{prop:multiscaling}).  The
        survivor ensemble produces this through the minimum-degree cutoff, and at leading order it
        is the whole bend.  The ratio curve is a leading-order property; $\beta$ is not.
  \item Open: $\beta_{\rm ad}$ on light tails, a nonlinear one-firm response ($\le0.12$); and
        $\alpha\lesssim2.2$, where a sparse, non-Gaussian cavity is needed.
\end{enumerate}

\end{document}
